\documentclass{article}
\usepackage{amsmath,amsthm,amssymb}
\usepackage[margin=1in]{geometry}

\newcommand{\degr}{\operatorname{degr}}

\title{\(q,t\)-Catalan Computer-Assisted Proofs 2024}
\author{}
\date{}

\begin{document}
\maketitle

\section{Purpose}

This item records the computer-assisted verification used for Lemma 2 and
Lemma 3 of Section 9 of Graham Hawkes's paper
\emph{A conjectured formula for the rational \(q,t\)-Catalan polynomial}
(Annals of Combinatorics 28, 749--795, 2024; arXiv:2208.00577).

The source computation is
\[
  \texttt{Conjectures-and-Computations/qt-catalan/qt-assisted.py}.
\]
The curated script
\[
  \texttt{code/qt\_assisted\_2024.py}
\]
is a reproducibility wrapper around the same finite checks.  It keeps the
source conventions and default cutoff \(d^*=20\), but replaces top-level script
side effects by a command-line entry point and stable summary output.

\section{Objects}

The computation uses position coordinates for \(m\)-Dyck paths.  A path is a
finite sequence
\[
  A=(a_0,\ldots,a_\ell),\qquad a_0=0,\qquad
  0\le a_{i+1}\le a_i+m.
\]
For integers \(a,b\), define
\[
  \alpha(a,b;m)=
  \begin{cases}
    \min(b-a,m), & a\le b,\\
    \min(a-b-1,m), & a>b,
  \end{cases}
  \qquad
  \alpha_0(a;m)=\max(0,a-m).
\]
The script generates paths together with the degree statistic obtained
incrementally from these \(\alpha\)-terms.  When appending a new final entry
\(j\) to a prefix \(x=(a_0,\ldots,a_{r-1})\), the degree increment is
\[
  \sum_{k=1}^{r-1}\alpha(a_k,j;m)-\alpha_0(j;m).
\]

The finite cutoff is \(d^*=20\).  For each \(1\le m\le20\), the script sets
\[
  \ell^*(m)=\left\lceil\frac{20}{m}+1.001\right\rceil
\]
and generates every path of length at most \(\ell^*(m)+1\) whose generated
degree is at most \(20\).  These are exactly the paths needed by the source
computation for the two Section 9 lemmas.

\section{Lemma 2 Check}

The source code's Lemma 2 check is a finite string-length bound.  It uses the
paper's point, pair, right, and lowest operations on position-coordinate paths.
For a path \(A\), let \(b=\operatorname{lowest}(A)\) be the terminal path
obtained by repeatedly applying the right map.  The computation also uses the
height
\[
  h(A)=(g-1)\ell+j-\sum_i a_i,
\]
where \(g\) is the largest entry of \(A\) and \(j\) is the last index at which
that largest entry is attained.

Only the boundary case needs checking: paths of length \(\ell^*(m)+1\) with
second entry \(0\).  For each such path, the script verifies
\[
  \operatorname{area}(b)\le
  m\binom{\ell^*(m)+1}{2}-h(A)-\degr(A).
\]
All shorter paths, and the length-boundary paths outside this maximal class,
are treated as automatic cases just as in the source script.

\section{Lemma 3 Check}

The Lemma 3 computation is a finite monomial identity checked separately for
each pair \((m,\ell)\) appearing in the generated records.  Put
\[
  M=m\binom{\ell+1}{2}.
\]
For every generated path \(A\) of fixed \(m\) and length \(\ell+1\), the left
side contributes the monomial record
\[
  (\operatorname{area}(A),\,M-\degr(A)).
\]
The right side is built from the subfamily with second entry \(0\).  If
\(a=\operatorname{area}(A)\) and \(d=\degr(A)\), this subfamily contributes a
positive interval
\[
  (j,M-d),\qquad a\le j\le M-a-d,
\]
when \(a\le M-a-d\), and a negative correction interval
\[
  (j,M-d),\qquad M-a-d+1\le j<a,
\]
when \(M-a-d<a\).  The check is
\[
  \text{positive interval multiset}
  =
  \text{all path monomials}+\text{negative correction multiset}.
\]
The script sorts both multisets by the source ordering and compares them
exactly.

\section{Reproducible Run}

Run the curated checker from the repository root by executing
\[
  \texttt{python items/qt\_catalan\_computer\_assisted\_proofs\_2024/code/qt\_assisted\_2024.py}
\]
from inside the \(\texttt{Combinatorics}\) directory, or equivalently by giving
the full path to the script.

The default run performed for this item generated \(5{,}692{,}942\) records.
Both checks passed:
\[
  \texttt{lemma2\_status: PASS},\qquad
  \texttt{lemma3\_status: PASS}.
\]
The Lemma 3 comparison checked 106 nontrivial \((m,\ell)\)-layers.  The final
layer in the default source ordering was \(m=1\), \(\ell=22\), with 607,712
records.  The measured run time in this environment was about 57 seconds.

\section{Status}

This is proof-supporting computation, not experimental evidence for a new
conjecture.  The mathematical proof in the 2024 paper reduces the remaining
parts of Lemma 2 and Lemma 3 to the finite checks above.  The curated code is
intended to make those checks reproducible from the archived source script.

\end{document}
